% formal/active-orbit-facet-coverage.tex -- input by formal/main.tex
%
% Purpose: isolate the exact fixed-normal support-number statement behind the
% heuristic that every facet of a local maximizer must be used by a minimizing
% orbit.  The full dual-vertex analogue requires a complete first-order branch
% model and is stated only conditionally below.
%
% Mathematical status (2026-07-22): agent-proved, not yet reviewed by Jörn.
% The proof depends on the HK2017 capacity formula in the convention audited in
% formal/hk2017-qp-conventions.tex and on the standard support-number first
% variation of volume.  No finite HKO active-orbit catalogue is asserted.

\subsection{Active-orbit balance and facet coverage}
\label{sec:active-orbit-facet-coverage}

\begin{unverified}
\paragraph{Status and scope.}
The results in this subsection are agent-proved and await Jörn's mathematical
review.  The unconditional statement is in a fixed-normal support-number
chamber.  It applies in particular to the support-number subspace of a stable
dual-vertex chart, but it is not a theorem about arbitrary changes of the
facet normals.  The latter requires the complete smooth-envelope hypothesis
in Lemma~\ref{lem:facet-coverage-abstract-balance}.

\paragraph{Research relevance.}
This subsection supplies a cheap necessary-condition mechanism for the search
for local maxima of \(\operatorname{sys}\).  In a stable fixed-normal chart,
failure of facet coverage is not merely suggestive: moving an uncovered facet
inward increases \(\operatorname{sys}\) to first order.  The stronger convex
balance in Theorem~\ref{thm:facet-coverage-active-weight-balance} can exclude a
putative local maximum even when every facet is covered separately.  These
tests can therefore reject candidates or explain improving directions before a
second-order analysis or a local optimization run.

The result does not by itself turn a numerical orbit list into such a test.
One must first know that the list represents every active limiting HK datum in
the chart, or have exact gap/dominance arguments for the omitted data.  This is
especially important at singular or highly symmetric examples.  The theorem's
direct thesis value is thus a reusable necessary condition and a precise
interface between orbit data and local-maximality claims, rather than a new
classification of local maxima.

Let $n_1,\ldots,n_F\in S^3$ be fixed labelled outward unit normals and put
\[
  K(h)=\{x\in\mathbb R^4:\langle n_i,x\rangle\leq h_i,
       \ i=1,\ldots,F\}.
\]
We work on an open support-number chamber $\mathcal H\subset\mathbb R^F$ on
which $0\in\operatorname{int}K(h)$, the body is bounded, every listed
inequality defines a facet, and the labelled face combinatorics is constant.
Thus all $h_i$ are positive and the volume $V(h)$ is smooth.  Write $S_i(h)$
for the three-dimensional volume of the $i$th facet.  Requiring such a chamber
to be open in all of $\mathbb R^F$ deliberately excludes nonsimple
support-wall points; no claim below covers a nonsimple local maximum.

For a permutation $\sigma\in S_F$, define the normals-form quadratic
expression in the project active-word orientation by
\[
  \mathcal Q_\sigma(b)=
  \sum_{1\leq j<k\leq F}b_{\sigma(j)}b_{\sigma(k)}
  \omega_0(n_{\sigma(j)},n_{\sigma(k)}).
\]
This has the orientation audited in
Subsection~\ref{subsec:hk2017-active-word-qp-convention}, but uses the fixed
unit normals, not the dual vertices $a_i=n_i/h_i$.  The candidate family is
closed under word reversal, so this orientation change does not change the
scalar capacity envelope.  Define the compact
normal-closure simplex
\[
  \Delta_n=
  \left\{b\in\mathbb R^F:
    b_i\geq0,\quad \sum_i b_i=1,\quad \sum_i b_i n_i=0
  \right\}.
\]
Zero coordinates of $b$ represent facets omitted from the active word.  For
$b\in\Delta_n$ with $\mathcal Q_\sigma(b)>0$, set
\begin{equation}\label{eq:facet-coverage-action-branch}
  A_{\sigma,b}(h)=\frac{(b\cdot h)^2}{2\mathcal Q_\sigma(b)}.
\end{equation}
The set of \emph{active HK data} at $h_0\in\mathcal H$ is
\begin{equation}\label{eq:facet-coverage-active-data}
  \mathcal A(h_0)=
  \{(\sigma,b):\mathcal Q_\sigma(b)>0,
      \ A_{\sigma,b}(h_0)=c_{\mathrm{EHZ}}(K(h_0))\}.
\end{equation}
For $(\sigma,b)\in\mathcal A(h_0)$, put
\begin{equation}\label{eq:facet-coverage-normalized-weight}
  \alpha=\frac{b}{b\cdot h_0}.
\end{equation}
Then $\alpha_i\geq0$, $\sum_i\alpha_i n_i=0$, and
$\alpha\cdot h_0=1$.

\begin{lemma}[Attaining HK data and facet dwell]
\label{lem:facet-coverage-orbit-realization}
Every datum in $\mathcal A(h_0)$ determines, after deleting zero entries, an
action-minimizing simple generalized characteristic whose project active
traversal order is $\sigma$.  Equivalently, the same datum is indexed by the
reversed word in HK2017's displayed convention.  Its dwell-time fraction
associated with facet $i$ is
\begin{equation}\label{eq:facet-coverage-dwell-fraction}
  T_i=\alpha_i h_{0,i}.
\end{equation}
Thus $\alpha_i>0$ if and only if this characteristic has a positive
$i$-labelled velocity interval, and such an interval spends positive time on
facet $F_i$.  The corresponding segment may lie in a lower-dimensional face
of $F_i$; relative-interior visitation is not asserted.  When $\alpha_i=0$,
another labelled segment may still lie in an intersection contained in $F_i$.
\end{lemma}

\begin{proof}
The relations $\alpha\cdot h_0=1$ and
$\sum_i\alpha_i n_i=0$ give $T_i\geq0$, $\sum_iT_i=1$, and, for the facet
velocities $p_i=2J_0n_i/h_{0,i}$,
\[
  \sum_iT_i p_i=2J_0\sum_i\alpha_i n_i=0.
\]
Place the constant velocity equal to $p_{\sigma(k)}$ on an interval of
length $T_{\sigma(k)}$, in the project active traversal order
$k=1,\ldots,F$.  The resulting piecewise-affine path closes.  In the project
convention, its unnormalized dual-action integral has the positive sign
\[
  \int\langle-J_0\dot z,z\rangle
  =4\sum_{1\leq j<k\leq F}
       \alpha_{\sigma(j)}\alpha_{\sigma(k)}
       \omega_0(n_{\sigma(j)},n_{\sigma(k)})
  =4\mathcal Q_\sigma(\alpha)>0.
\]
The reversed word $\rho(k)=\sigma(F+1-k)$ enters only when translating this
datum to HK2017's displayed word convention.  Rescaling the positive-action
loop to the dual-action normalization used in the proof of the Haim--Kislev
formula gives dual value $A_{\sigma,b}(h_0)/2$.  Because the datum is active,
this is the minimum dual value.  The dual-action correspondence therefore
yields an action-minimizing generalized characteristic with these velocity
intervals.

On an interval with $T_i>0$, its velocity is a positive multiple of
$J_0n_i$.  The generalized characteristic inclusion then has
$n_i$ in the normal cone of $K(h_0)$ along that interval, which implies
$\langle n_i,x\rangle=h_{0,i}$ there.  Hence the curve lies in $F_i$ for a
positive time.  Conversely $T_i=0$ removes the $i$-labelled interval, without
excluding incidental contact with $F_i$ during another labelled interval.
Since $h_{0,i}>0$, \eqref{eq:facet-coverage-dwell-fraction} proves the stated
equivalence.
\end{proof}

We therefore say that an active datum or its characteristic \emph{uses} the
$i$th facet when $\alpha_i>0$, meaning positive $i$-labelled dwell.  This is
positive HK weight, not mere occurrence of the label in a padded word or
incidental contact along a lower-dimensional face.

\begin{lemma}[Exact support-number envelope and derivative]
\label{lem:facet-coverage-support-envelope}
For every $h\in\mathcal H$,
\begin{equation}\label{eq:facet-coverage-capacity-envelope}
  c_{\mathrm{EHZ}}(K(h))
  =\min_{\substack{\sigma, b\in\Delta_n\\\mathcal Q_\sigma(b)>0}}
       A_{\sigma,b}(h).
\end{equation}
At $h_0$, the one-sided directional derivative exists in every
$v\in\mathbb R^F$ and is
\begin{equation}\label{eq:facet-coverage-capacity-derivative}
  Dc_{\mathrm{EHZ}}(h_0;v)
  =2c_0\min_{\alpha\in\mathcal W(h_0)}\alpha\cdot v,
  \qquad c_0=c_{\mathrm{EHZ}}(K(h_0)),
\end{equation}
where $\mathcal W(h_0)$ is the compact set of normalized weights
\eqref{eq:facet-coverage-normalized-weight} coming from
$\mathcal A(h_0)$.
\end{lemma}

\begin{proof}
The normals/heights form of the Haim--Kislev formula maximizes
$\mathcal Q_\sigma(\beta)$ subject to
$\beta\geq0$, $\sum_i\beta_i n_i=0$, and $\beta\cdot h=1$.
Writing $\beta=b/(b\cdot h)$ with $b\in\Delta_n$ turns its reciprocal into
\eqref{eq:facet-coverage-capacity-envelope}.

There is a minor compactness point because the displayed minimum excludes
$\mathcal Q_\sigma(b)\leq0$.  On a sufficiently small neighborhood of $h_0$ there is
$m>0$ with $b\cdot h\geq m$ for every $b\in\Delta_n$.  A fixed feasible
positive-$Q$ datum gives a uniform upper bound $M$ for the minimum.  Every
datum with $A_{\sigma,b}(h)\leq M$ then satisfies
$\mathcal Q_\sigma(b)\geq m^2/(2M)$.  Thus all local minimizers lie in a compact
subset on which the branch functions are smooth.  In particular the active
set is compact.

The directional derivative is now the elementary compact-envelope formula
\[
  Dc_{\mathrm{EHZ}}(h_0;v)
  =\min_{(\sigma,b)\in\mathcal A(h_0)}
    \frac{(b\cdot h_0)(b\cdot v)}{\mathcal Q_\sigma(b)}.
\]
For completeness, the upper bound follows by keeping any active datum fixed.
For the lower bound, choose minimizing data at $h_0+tv$, pass to a convergent
subsequence in the compact set above, and use that the limit is active at
$h_0$; the nonnegative base-point gap can only increase the difference
quotient.  Finally, activity gives
$(b\cdot h_0)^2/(2\mathcal Q_\sigma(b))=c_0$, so the last display is exactly
\eqref{eq:facet-coverage-capacity-derivative}.
\end{proof}

This proof also settles the right-active-germ issue inside this chart.  The
feasible set $\Delta_n$ and every $\mathcal Q_\sigma$ are independent of $h$.  Hence a
nearby minimizing germ has a base-point limit in $\mathcal A(h_0)$; a
base-invalid positive-$Q$ germ cannot appear.  Singular KKT matrices do not
enter the argument.  Zero coordinates are retained: if $b_i(t)>0$ tends to
$b_i(0)=0$, the compact-envelope derivative still uses the partial
$h$-derivative at $b(0)$, whose $i$th component is zero.  A finite numerical
orbit list may replace $\mathcal A(h_0)$ only after proving that it contains
all active limits (or proving exact gap/dominance certificates for what it
discards).

The volume derivative in this chamber is
\begin{equation}\label{eq:facet-coverage-volume-derivative}
  DV(h_0)[v]=\sum_{i=1}^F S_i v_i.
\end{equation}
Indeed, increasing $h_i$ moves the $i$th facet outward at unit normal speed.
Translation and scaling of the body also give the familiar identities
\begin{equation}\label{eq:facet-coverage-minkowski-identities}
  \sum_i S_i n_i=0,
  \qquad
  \sum_i S_i h_{0,i}=4V_0,
  \qquad V_0=V(h_0).
\end{equation}

\begin{theorem}[Active-weight balance at a local maximum]
\label{thm:facet-coverage-active-weight-balance}
If $h_0$ is a local maximum of
$\operatorname{sys}(K(h))=c_{\mathrm{EHZ}}(K(h))^2/(2V(h))$ in
$\mathcal H$, then
\begin{equation}\label{eq:facet-coverage-convex-balance}
  \left(\frac{S_1}{4V_0},\ldots,\frac{S_F}{4V_0}\right)
  \in\operatorname{conv}\mathcal W(h_0).
\end{equation}
Equivalently, there are finitely many active normalized weights
$\alpha^{(k)}$ and coefficients $\lambda_k\geq0$, $\sum_k\lambda_k=1$, such
that
\begin{equation}\label{eq:facet-coverage-coordinate-balance}
  \sum_k\lambda_k\alpha_i^{(k)}=\frac{S_i}{4V_0}
  \quad(i=1,\ldots,F).
\end{equation}
One may take at most $F+1$ active data.
\end{theorem}

\begin{proof}
Combining Lemma~\ref{lem:facet-coverage-support-envelope} with
\eqref{eq:facet-coverage-volume-derivative} gives
\begin{equation}\label{eq:facet-coverage-sys-derivative}
  D\operatorname{sys}(h_0;v)
  =\operatorname{sys}(h_0)
   \min_{\alpha\in\mathcal W(h_0)}
   \left(4\alpha-\frac{S}{V_0}\right)\!\cdot v,
\end{equation}
where $S=(S_1,\ldots,S_F)$.  At a local maximum this is nonpositive for every
$v$.  If zero were not in the convex hull of
$\{4\alpha-S/V_0:\alpha\in\mathcal W(h_0)\}$, strict separation in
$\mathbb R^F$ would give a $v$ on which every functional in the braces is
positive, contradicting \eqref{eq:facet-coverage-sys-derivative}.  This proves
\eqref{eq:facet-coverage-convex-balance}; the finite form follows from
Carath\'eodory's theorem.
\end{proof}

\begin{corollary}[Facet coverage, with a quantitative bound]
\label{cor:facet-coverage-active-orbit}
Under the hypotheses of
Theorem~\ref{thm:facet-coverage-active-weight-balance}, every facet is used by
at least one active minimizing HK datum.  More precisely, for every $i$,
\begin{equation}\label{eq:facet-coverage-quantitative}
  \max_{\alpha\in\mathcal W(h_0)}\alpha_i
  \geq \frac{S_i}{4V_0}>0.
\end{equation}
Equivalently, when only the $i$th support number moves inward,
\begin{equation}\label{eq:facet-coverage-inward-capacity-drop}
  Dc_{\mathrm{EHZ}}(h_0;-e_i)
  =-2c_0\max_{\alpha\in\mathcal W(h_0)}\alpha_i
  \leq-\frac{c_0S_i}{2V_0}.
\end{equation}
If every active datum were invisible to facet $i$, meaning
$\alpha_i=0$, capacity would be stationary to first order in this inward
direction while volume would have derivative $-S_i$; then
$D\operatorname{sys}(h_0;-e_i)=\operatorname{sys}(h_0)S_i/V_0>0$.
\end{corollary}

\begin{proof}
Each $S_i$ is positive because every listed inequality is facet-defining.
Taking the $i$th coordinate in
\eqref{eq:facet-coverage-coordinate-balance} proves
\eqref{eq:facet-coverage-quantitative}.  Formula
\eqref{eq:facet-coverage-inward-capacity-drop} follows directly from
\eqref{eq:facet-coverage-capacity-derivative}.
\end{proof}

The quantifier in Corollary~\ref{cor:facet-coverage-active-orbit} is
facetwise: for every facet there exists a minimizing orbit that uses it.  The
orbit may depend on the facet.  The argument does not produce one minimizing
orbit that uses all facets, and no such stronger condition is needed here.
Jörn has confirmed that the intended research question uses this facetwise
quantifier.

\paragraph{Quotient interpretation.}
Within the fixed-normal chart, translations and positive scalings have tangent
space
\[
  T=\{v\in\mathbb R^F:v_i=\langle n_i,y\rangle+\rho h_{0,i}
      \text{ for some }y\in\mathbb R^4,\ \rho\in\mathbb R\}.
\]
The rows $4\alpha-S/V_0$ in
\eqref{eq:facet-coverage-sys-derivative} annihilate $T$ by closure,
normalization, and \eqref{eq:facet-coverage-minkowski-identities}.  Thus
\eqref{eq:facet-coverage-convex-balance} is equivalently a convex-hull balance
in $(\mathbb R^F/T)^*$.  General linear symplectic directions change the
normals and are not directions in this chart.  A local maximum in the full
dual-vertex chart is nevertheless a local maximum on the support-number
subchart because
\begin{equation}\label{eq:facet-coverage-support-to-row}
  a_i=\frac{n_i}{h_i},
  \qquad
  \dot a_i=-\frac{v_i}{h_i}a_i.
\end{equation}
Consequently the facet-coverage corollary is necessary for a full-chart local
maximum as well, although it is only the radial-row projection of any
full-chart balance condition.

\begin{lemma}[Conditional balance in an arbitrary chart]
\label{lem:facet-coverage-abstract-balance}
Let $X$ be a finite-dimensional chart, $T\subset X$ a symmetry tangent space,
and $S$ a complement.  Suppose a function $f$ has, uniformly for $x\to0$ in
$S$, a complete first-order representation
\[
  f(x)=f(0)+\min_{\gamma\in\Gamma}r_\gamma(x)+o(\|x\|),
\]
where $\Gamma$ is compact, the $r_\gamma\in S^*$ depend continuously on
$\gamma$, and every right-active, singular, and zero-weight germ is included.
If $0$ is a local maximum on the slice, then
\[
  0\in\operatorname{conv}\{r_\gamma:\gamma\in\Gamma\}\subset S^*.
\]
\end{lemma}

\begin{proof}
Otherwise strict separation gives $v\in S$ and $\varepsilon>0$ with
$r_\gamma(v)\geq\varepsilon$ for every $\gamma$.  The uniform expansion then
gives $f(tv)>f(0)$ for all sufficiently small $t>0$, a contradiction.
\end{proof}

For a complete finite family of smooth optimizing action branches in the
dual-vertex chart, Lemma~\ref{lem:facet-coverage-abstract-balance} is the usual
quotient-gradient convex-hull condition.  Base-point KKT rows alone do not
establish its hypothesis: beta-zero appearances, singular optimizer sets, and
discarded words require a complete right-active model or exact gap/dominance
certificates.  In particular, the selected HKO feasible sections in
Section~\ref{sec:hko-feasible-section-upper-branches} are touching upper
functions used for a sufficient certificate; they are deliberately not a
complete active-orbit catalogue and should not be used to claim the necessary
condition in this lemma.

\paragraph{Sufficiency boundary.}
Facet coverage is only a coordinatewise consequence of
\eqref{eq:facet-coverage-convex-balance}.  It does not imply convex balance and
therefore does not by itself certify local maximality.  Indeed, an abstract
lower-envelope family can contain a row that sees every coordinate while its
row convex hull is strictly separated from zero, producing a common increasing
direction.  This is a logical counterexample to coverage as a first-order
certificate; no actual EHZ support-chart polytope with coverage but without
local maximality is claimed here.  The balance condition itself is not
sufficient when zero lies on the boundary: first-order flat directions require
higher order control (the one-function example $f(x)=x^2$ already has zero
gradient at a local minimum).

A stronger first-order hypothesis can be sufficient.  On a complement to
$T$, if zero lies in the interior of the convex hull of the active rows
$4\alpha-S/V_0$, then compactness gives $\eta>0$ such that every unit slice
direction has some active row at most $-\eta$.  The smooth upper functions
\[
  h\longmapsto \frac{A_{\sigma,b}(h)^2}{2V(h)}
\]
touch $\operatorname{sys}$ at $h_0$ and dominate it nearby.  Their uniform
first-order expansions then prove a strict local maximum on the slice.  This
is a strict convex-hull/spanning certificate, not facet coverage.  Boundary
balance needs a separate second-order or exact upper-function argument.
\end{unverified}
